\documentclass[10pt,a4paper]{article}
\usepackage[a4paper,top=20mm,bottom=22mm,left=16mm,right=16mm]{geometry}
\usepackage{fontspec}
\setmainfont{Noto Sans CJK KR}
\setsansfont{Noto Sans CJK KR}
\setmonofont{Noto Sans Mono CJK KR}[Scale=0.78]
\usepackage{unicode-math}
\setmathfont{Latin Modern Math}
\usepackage{listings}
\usepackage{xcolor}
\usepackage{amsmath}
\usepackage{booktabs}
\usepackage{array}
\usepackage{fancyhdr}
\usepackage{titlesec}
\usepackage[hidelinks]{hyperref}
\usepackage{enumitem}
\usepackage{tcolorbox}
\tcbuselibrary{skins,breakable}

\definecolor{codebg}{HTML}{F7F7F4}
\definecolor{kw}{HTML}{1F5FA8}
\definecolor{cmt}{HTML}{6B7280}
\definecolor{str}{HTML}{2E7D4F}
\definecolor{num}{HTML}{C1382B}
\definecolor{rule}{HTML}{9CA3AF}
\definecolor{accent}{HTML}{1F5FA8}

\lstset{
  basicstyle=\ttfamily\scriptsize,
  numbers=left, numberstyle=\ttfamily\tiny\color{rule},
  numbersep=6pt,
  keywordstyle=\color{kw}\bfseries,
  commentstyle=\color{cmt}\itshape,
  stringstyle=\color{str},
  showstringspaces=false,
  breaklines=true, breakatwhitespace=false,
  breakindent=10pt,
  columns=fullflexible, keepspaces=true,
  backgroundcolor=\color{codebg},
  frame=leftline, framerule=1.1pt, rulecolor=\color{rule},
  framexleftmargin=4pt, xleftmargin=14pt,
  aboveskip=5pt, belowskip=5pt,
  language=Python,
  literate={~}{{\textasciitilde}}1
}

\newcommand{\note}[1]{%
  \begin{tcolorbox}[enhanced,breakable,colback=white,colframe=accent,
    boxrule=0pt,leftrule=1.6pt,arc=0pt,outer arc=0pt,
    left=6pt,right=2pt,top=3pt,bottom=3pt,boxsep=0pt,
    before skip=5pt, after skip=8pt]
  #1
  \end{tcolorbox}}

\titleformat{\section}{\Large\bfseries\color{accent}}{\thesection}{0.6em}{}
\titleformat{\subsection}{\large\bfseries}{\thesubsection}{0.6em}{}
\titleformat{\subsubsection}{\normalsize\bfseries}{\thesubsubsection}{0.6em}{}
\titlespacing*{\section}{0pt}{16pt}{7pt}
\titlespacing*{\subsection}{0pt}{11pt}{5pt}

\pagestyle{fancy}\fancyhf{}
\fancyhead[L]{\footnotesize\color{cmt} optical-metrology-fitter — 물리 해설판}
\fancyhead[R]{\footnotesize\color{cmt}\thepage}
\renewcommand{\headrulewidth}{0.3pt}

\tolerance=3500
\emergencystretch=3em
\hbadness=10000
\setlength{\parindent}{0pt}
\setlength{\parskip}{4pt}
\newcommand{\nt}{\tilde n}
\newcommand{\Imag}{\operatorname{Im}}
\newcommand{\Real}{\operatorname{Re}}

\begin{document}

\begin{titlepage}
\vspace*{2.2cm}
{\Huge\bfseries\color{accent} optical-metrology-fitter}\\[6pt]
{\LARGE 코드 전문과 물리 해설}\\[22pt]
{\large 전달행렬법 기반 박막 광학 계측 —\\ 순방향 모델, 역문제, 그리고 불확도 예산}\\[30pt]
\rule{\textwidth}{0.7pt}\\[10pt]
{\normalsize
\textbf{규약} \quad $e^{-i\omega t}$, \quad $\nt = n + ik$ ($k\ge 0$), \quad 순방향파 $e^{+ik_z z}$\\[4pt]
\textbf{핵심 결과} \quad 알고리즘 잡음 바닥 $\approx 4$ pm \quad vs \quad 현실적 정확도 $\approx 1.3$ nm \quad (격차 $\approx 300\times$)\\[4pt]
\textbf{검증} \quad 34/34 통과, Byrnes 레퍼런스 구현체와 $10^{-14}$ 수준 일치
}\\[8pt]
\rule{\textwidth}{0.7pt}

\vfill
{\footnotesize\color{cmt}
이 문서는 저장소의 모든 소스 코드를, 각 코드 덩어리 바로 옆(또는 아래)에 대응하는 물리적 근거를 붙여 재구성한 것입니다.
Part I은 코드와 독립적인 물리 유도, Part II는 주석 달린 소스, Part III는 축약 없는 전체 원본입니다.\\[6pt]
2026년 8월}
\end{titlepage}

\tableofcontents
\newpage

%==============================================================
\section{읽기 전에 — 이 문서의 구조}

이 프로젝트가 답하려는 질문은 하나입니다.

\begin{center}
\emph{반사율 $R$을 측정했을 때, 박막 두께 $d$를 얼마나 믿을 수 있는가?}
\end{center}

인과 사슬은 다음과 같고, 문서 전체가 이 순서를 따릅니다.

\vspace{4pt}
\begin{center}
\fbox{Maxwell 경계조건} $\to$ \fbox{Fresnel $r,t$} $\to$ \fbox{다중반사 합} $\to$ \fbox{$R(\lambda,\theta;d)$} $\to$ \fbox{$d$ 역추정} $\to$ \fbox{불확도}
\end{center}
\vspace{4pt}

핵심 명제를 미리 못 박아 둡니다. \textbf{박막 두께는 직접 관측량이 아닙니다.} 관측량은 왕복 위상 $2\delta$ 하나뿐이고 두께는 거기서 \emph{추론}됩니다. 이 프로젝트의 모든 어려움 — 프린지 차수 축퇴, $n\cdot d$ 축퇴, 300배 격차 — 이 전부 이 한 문장에서 파생됩니다.

\subsection*{문서 구성}

\begin{itemize}[leftmargin=14pt,itemsep=2pt]
\item \textbf{Part I (§2--§7)} — 코드와 무관하게 성립하는 물리. Maxwell 경계조건에서 Cramér--Rao 하한까지.
\item \textbf{Part II (§8--§12)} — 소스 코드 전문. 코드 덩어리마다 바로 아래에 파란 세로선으로 표시된 물리 해설을 붙였습니다. 표시된 줄 번호는 각 모듈 원본 파일의 실제 줄 번호가 아니라 덩어리 내부의 상대 번호입니다.
\item \textbf{Part III (§13)} — 축약 없는 전체 원본 소스. Part II에서 지면상 접었던 docstring까지 전부 포함합니다.
\end{itemize}

\note{\textbf{Part II의 표기 약속.} 코드 덩어리는 회색 배경에 파란 줄번호로, 물리 해설은 파란 세로선 상자로 구분됩니다. Part II의 코드 덩어리에서는 여러 줄 docstring을 첫 줄로 접었습니다 — 그 내용이 곧 옆의 해설이기 때문입니다. 접지 않은 완전한 원본은 Part III에 있습니다.}

%==============================================================
\part*{\color{accent}Part I — 물리}
\addcontentsline{toc}{section}{Part I — 물리}

\section{Fresnel 계수 — 사실은 임피던스 정합}

\subsection{왜 경계조건에서 출발하는가}

자유전하·자유전류가 없는 계면에 Maxwell 방정식의 적분형을 적용하면 단 두 개의 조건이 남습니다.

\begin{equation}
\hat n \times (\mathbf{E}_2 - \mathbf{E}_1) = 0, \qquad
\hat n \times (\mathbf{H}_2 - \mathbf{H}_1) = 0
\end{equation}

즉 \textbf{접선 성분 $E$, $H$의 연속}입니다. 이것이 전부이고, 나머지는 순수한 대수입니다. 여기서 유의할 점은 $\mathbf{D}$와 $\mathbf{B}$의 법선 성분 연속은 위 두 조건의 종속 결과라는 것입니다 — 독립 조건이 아니므로 따로 쓸 필요가 없습니다.

\subsection{\texorpdfstring{s편광 (TE): $\mathbf{E}\parallel\hat y$}{s편광 (TE)}}

$E_y$ 연속으로부터
\begin{equation}
1 + r = t
\end{equation}
$\mathbf{H} = (\mathbf{k}\times\mathbf{E})/(\omega\mu_0)$이므로 $H_x \propto k_z E_y$이고, $H_x$ 연속으로부터
\begin{equation}
k_{z1}(1-r) = k_{z2}\,t
\end{equation}
두 식을 연립하면
\begin{equation}
\boxed{\;r_s = \frac{k_{z1}-k_{z2}}{k_{z1}+k_{z2}}, \qquad
t_s = \frac{2k_{z1}}{k_{z1}+k_{z2}}\;}
\end{equation}

\note{\textbf{직관.} 이 식은 전송선 임피던스 정합 공식 $\Gamma=(Z_2-Z_1)/(Z_2+Z_1)$과 글자 그대로 같은 형태입니다. $k_z$가 광학 어드미턴스 역할을 합니다. 임피던스가 맞으면 반사가 없다는 전기공학의 직관이 그대로 광학에 이식됩니다.}

\subsection{\texorpdfstring{p편광 (TM): $\mathbf{H}\parallel\hat y$}{p편광 (TM)}}

$H_y$ 연속으로부터 $1+r = t_H$. 그리고 $\nabla\times\mathbf{H} = -i\omega\varepsilon\mathbf{E}$에서 $E_x \propto (k_z/\varepsilon)H_y$이므로
\begin{equation}
\frac{k_{z1}}{\varepsilon_1}(1-r) = \frac{k_{z2}}{\varepsilon_2}t_H
\end{equation}
따라서
\begin{equation}
\boxed{\;r_p = \frac{\varepsilon_2 k_{z1}-\varepsilon_1 k_{z2}}{\varepsilon_2 k_{z1}+\varepsilon_1 k_{z2}}
= \frac{\nt_2^2 k_{z1}-\nt_1^2 k_{z2}}{\nt_2^2 k_{z1}+\nt_1^2 k_{z2}}\;}
\end{equation}

코드의 $t_p = 2\nt_1\nt_2 k_{z1}/(\nt_2^2k_{z1}+\nt_1^2k_{z2})$는 위 결과를 $H$ 진폭이 아니라 \textbf{$E$ 진폭}으로 환산한 것입니다 ($|E| = |H|/\nt$).

\subsection{s와 p의 통합 — 광학 어드미턴스}

\begin{equation}
r = \frac{Y_1-Y_2}{Y_1+Y_2}, \qquad
Y_s \propto k_z, \qquad Y_p \propto \frac{\nt^2}{k_z}
\end{equation}

s와 p는 다른 물리가 아니라 \textbf{다른 어드미턴스}입니다. 박막 설계 문헌 전체가 이 형식으로 쓰여 있습니다.

\note{\textbf{함정 — $r_p$의 부호는 규약입니다.} $Y_p$ 정의로 계산하면 이 프로젝트의 $r_p$와 부호가 반대로 나옵니다. 이는 반사 후 $\mathbf{E}_p$의 양의 방향을 어떻게 잡느냐의 순수한 규약 문제입니다. $|r_p|^2$는 규약과 무관하므로 $R_p$는 안전하지만, 서로 다른 출처의 $r_p$ 정의를 한 계산에 섞으면 반드시 사고가 납니다. 이 저장소는 Byrnes 규약($\theta_B$ 아래에서 $r_p>0$)을 일관되게 씁니다.}

\subsection{Brewster 각 — 쌍극자 복사 패턴}

$\theta_1+\theta_2 = 90°$일 때 $r_p = 0$이고, 이때 $\tan\theta_B = n_2/n_1$입니다. 물리적 이유는 다음과 같습니다.

매질 2에서 유도된 쌍극자는 투과파의 $\mathbf{E}_p$ 방향으로 진동합니다. p편광에서 이 방향은 반사파가 나가야 할 방향과 정확히 평행해집니다. 그런데 쌍극자 복사는 $\propto\sin^2\Theta$이므로 \textbf{자기 축 방향으로는 복사하지 않습니다}. 따라서 반사파가 아예 생성되지 않습니다.

\note{\textbf{흡수가 있으면 왜 0이 안 되는가.} $k\neq0$이면 쌍극자 응답에 위상 지연이 생겨 상쇄가 불완전해집니다. 실험 1에서 405\,nm의 Si는 $R_{p,\min}=5.2\times10^{-4}$인 반면, 무손실 BK7은 $4\times10^{-8}$입니다. \textbf{이 잔류량 자체가 $k$의 직접 측정량}이며, 엘립소미트리의 물리적 근거입니다.}

%==============================================================
\section{위상 규약과 분지 선택}

\subsection{세 선택은 독립이 아니다}

\begin{center}
\begin{tabular}{@{}ll@{}}
\toprule
선택 & 이 프로젝트 \\
\midrule
시간 의존성 & $e^{-i\omega t}$ \\
복소 굴절률 & $\nt = n+ik$, $k\ge0$ \\
순방향파 & $e^{+ik_z z}$ \\
\bottomrule
\end{tabular}
\end{center}

$+z$ 진행파의 진폭을 전개하면
\begin{equation}
e^{ik_z z} = \underbrace{e^{i\Real(k_z)z}}_{\text{진동}}\cdot\underbrace{e^{-\Imag(k_z)z}}_{\text{감쇠}}
\end{equation}

수동 매질에서 감쇠하려면 $\Imag(k_z)\ge0$이어야 합니다. \textbf{이것이 코드가 필요로 하는 유일한 분지 규칙입니다.}

반대 분지를 고르면 (i) 거리에 따라 증폭되는 비물리적 파가 되고, (ii) 두꺼운 흡수층에서 $e^{ik_zd}$가 오버플로를 일으켜 계산이 터집니다.

\note{\textbf{규약 혼용 경고.} 엘립소미트리 문헌의 상당수는 $e^{+i\omega t}$, $\nt=n-ik$를 씁니다. 두 규약 모두 옳지만 섞으면 반드시 틀립니다. 그래서 이 저장소는 \texttt{core.py} 최상단에 규약을 명시하고 이후 어디서도 다시 건드리지 않습니다.}

\subsection{스넬 법칙을 각도로 쓰지 않은 이유}

면내 병진 대칭성으로부터 면내 파수가 보존됩니다.
\begin{equation}
k_x = n_0 k_0 \sin\theta_0 = \text{const},\qquad
k_{z,j} = \sqrt{(\nt_j k_0)^2 - k_x^2}
\end{equation}

이것이 \textbf{스넬 법칙 그 자체}입니다. $\theta_j = \arcsin(\cdot)$로 쓰는 대신 $k_z$로 쓰면 다음이 자동으로 처리됩니다.

\begin{itemize}[leftmargin=14pt,itemsep=1pt]
\item 전반사 $\to$ $k_z$가 순허수
\item 흡수 매질 $\to$ $k_z$가 복소수
\item 소멸파(frustrated TIR) $\to$ 위와 동일
\end{itemize}

각도로 쓰면 복소 arcsin이 어느 리만 시트에 있는지를 매번 추적해야 합니다. \textbf{특수 케이스가 하나도 없는 코드}가 나온 이유가 이것입니다.

%==============================================================
\section{다중 반사 — Airy 급수와 전달행렬은 같은 것}

\subsection{경로 합}

박막 내부에서 빛이 무한히 왕복하며 매 왕복마다 위상 $2\delta$를 얻습니다.
\begin{equation}
\delta = k_{z1}d = \frac{2\pi}{\lambda}\nt_1 d\cos\theta_1
\end{equation}

전체 반사 진폭은 모든 경로의 코히런트 합입니다.
\begin{equation}
r = r_{01} + t_{01}r_{12}t_{10}e^{2i\delta} + t_{01}r_{12}r_{10}r_{12}t_{10}e^{4i\delta}+\cdots
\end{equation}
등비급수를 합하면
\begin{equation}
r = r_{01} + \frac{t_{01}t_{10}r_{12}e^{2i\delta}}{1-r_{10}r_{12}e^{2i\delta}}
\end{equation}

\subsection{Stokes 관계 — 시간 역전 대칭의 산물}

무손실 계면에서 시간 역전 불변성으로부터
\begin{equation}
t_{01}t_{10} = 1-r_{01}^2, \qquad r_{10} = -r_{01}
\end{equation}
이를 대입하면
\begin{equation}
\boxed{\;r = \frac{r_{01}+r_{12}e^{2i\delta}}{1+r_{01}r_{12}e^{2i\delta}}\;}\qquad\text{(Airy 공식)}
\end{equation}

\subsection{전달행렬은 이 급수를 미리 합해둔 것}

\begin{equation}
M = D_{01}\prod_{j=1}^{N-2}P_j D_{j,j+1},\quad
D_{ij}=\frac{1}{t_{ij}}\begin{pmatrix}1&r_{ij}\\r_{ij}&1\end{pmatrix},\quad
P_j=\begin{pmatrix}e^{-i\delta_j}&0\\0&e^{+i\delta_j}\end{pmatrix}
\end{equation}
\begin{equation}
r=\frac{M_{21}}{M_{11}},\qquad t=\frac{1}{M_{11}}
\end{equation}

$N=3$에 대해 직접 곱하면
\begin{equation}
M_{11}=\frac{e^{-i\delta}+r_{01}r_{12}e^{i\delta}}{t_{01}t_{12}},\qquad
M_{21}=\frac{r_{01}e^{-i\delta}+r_{12}e^{i\delta}}{t_{01}t_{12}}
\end{equation}
이므로 $r=M_{21}/M_{11}$이 정확히 Airy 공식으로 환원됩니다 (검증 B5, 오차 $2\times10^{-16}$).

\note{\textbf{행렬 형식의 이점.} 층이 $N$개여도 급수를 다시 유도할 필요 없이 행렬 곱만 늘리면 됩니다. 그리고 $D$(계면, 진폭 섞기)와 $P$(전파, 위상 축적)로 물리적 역할이 깔끔히 분리됩니다. 이 분리가 코드 구조에 그대로 반영되어 있습니다.}

\subsection{반파장층은 왜 투명한가}

$\delta=\pi$이면 $e^{2i\delta}=1$이므로
\begin{equation}
r = \frac{r_{01}+r_{12}}{1+r_{01}r_{12}} = r_{02}
\end{equation}
\textbf{막이 없는 것과 완전히 동일}합니다. 왕복 위상이 정확히 $2\pi$라 모든 경로가 동상으로 재결합해 계면 하나짜리 결과로 붕괴합니다. 검증 B6에서 $d_{1/2}=226.917$\,nm에 대해 오차 $1\times10^{-16}$로 확인했습니다.

%==============================================================
\section{투과율 — p편광에서 켤레가 필요한 이유}

$T$는 $|t|^2$가 아니라 \textbf{$z$방향 시간평균 Poynting 플럭스의 비}입니다.
\begin{equation}
S_z = \tfrac12\Real(\mathbf{E}\times\mathbf{H}^*)_z
\end{equation}

\textbf{s편광}: $H_x\propto k_z E_y$이므로 $S_z\propto\Real(k_z)|E|^2$, 따라서
\begin{equation}
T_s = |t|^2\,\frac{\Real(k_{zf})}{\Real(k_{z0})}
\end{equation}

\textbf{p편광}: $E_x\propto(k_z/\varepsilon)H_y$이고 $|H|=|\nt||E|$이므로
\begin{equation}
S_z \propto \Real\!\left(\frac{k_z}{\nt^2}\right)|\nt|^2|E|^2
= \frac{\Real(\nt^2k_z^*)}{|\nt|^2}|E|^2
\end{equation}
\begin{equation}
T_p = |t|^2\,\frac{\Real(\nt_f^2k_{zf}^*)/|\nt_f|^2}{\Real(\nt_0^2k_{z0}^*)/|\nt_0|^2}
\end{equation}

$\nt$가 실수면 두 식이 일치하지만 흡수 매질에서는 다릅니다. p에 s식을 쓰면 에너지 보존이 깨지고 오차는 각도와 $\Imag(\nt)$에 따라 커집니다 (검증 A3에서 검출).

\subsection{\texorpdfstring{반무한 흡수 기판에서 $T$의 의미}{반무한 흡수 기판에서 T의 의미}}

\note{\textbf{제가 실제로 틀렸던 지점입니다.} 처음 검증 항목을 세울 때 ``405\,nm에서 Si는 불투명하니 $T\to0$''을 기대했습니다. \textbf{테스트의 물리가 틀렸습니다.} 반무한 매질에는 두 번째 계면이 없으므로 $T$는 기판 \emph{앞면을 통과해 들어간} 플럭스입니다. 그 플럭스는 무한 경로에 걸쳐 흡수되지만 이미 들어갔습니다. 따라서 $R+T=1$이 정확히 성립하고, 이 부기 체계에서 $A=0$입니다. 불투명성은 \textbf{유한 슬랩}의 성질이며, 별도의 Beer--Lambert 검증(B13)으로 분리했습니다. 405\,nm에서 $1/\alpha = 124.3$\,nm입니다.}

%==============================================================
\section{물질 — 분산은 어디서 오는가}

\subsection{\texorpdfstring{SiO$_2$: Sellmeier는 감쇠 없는 Lorentz 진동자}{SiO2: Sellmeier는 감쇠 없는 Lorentz 진동자}}

\begin{equation}
\varepsilon(\omega) = 1+\sum_j\frac{f_j\omega_{pj}^2}{\omega_j^2-\omega^2-i\gamma_j\omega}
\end{equation}
투명 영역에서 $\gamma\to0$으로 두고 $\omega=2\pi c/\lambda$로 바꾸면
\begin{equation}
n^2 = 1+\sum_j\frac{B_j\lambda^2}{\lambda^2-C_j},\qquad
B_j=\frac{f_j\omega_{pj}^2}{\omega_j^2},\quad C_j=\lambda_j^2
\end{equation}

Malitson 계수의 물리적 정체는 다음과 같습니다.

\begin{center}
\begin{tabular}{@{}lll@{}}
\toprule
$\sqrt{C_j}$ & 에너지 & 정체 \\
\midrule
68.4 nm & 18.1 eV & 심준위 전자 전이 \\
116.2 nm & 10.7 eV & SiO$_2$ 밴드갭/엑시톤 (갭 $\sim$9 eV) \\
9.896 $\mu$m & 125 meV (1010 cm$^{-1}$) & Si--O--Si 신축 포논 \\
\bottomrule
\end{tabular}
\end{center}

즉 가시광 SiO$_2$의 분산은 \textbf{UV 전자 전이와 IR 포논 사이에 끼인 꼬리}입니다. 세 항은 임의로 맞춘 피팅 파라미터가 아닙니다.

\subsection{Si: 왜 Sellmeier가 안 되는가}

Si는 가시·UV에서 밴드간 흡수가 강해 $\gamma\to0$ 가정이 붕괴합니다. 저차 해석식이 존재하지 않으므로 측정 표를 보간해야 합니다. 임계점 구조는 다음과 같습니다.

\begin{itemize}[leftmargin=14pt,itemsep=1pt]
\item $E_1\approx3.4$ eV (365 nm): $\Lambda$ 방향 전이 $\to$ $n$이 6.03까지 상승
\item $E_2\approx4.25$ eV (292 nm): $X/\Sigma$ 근처, $\varepsilon_2$ 최대
\item 간접갭 1.12 eV (1100 nm): 포논 보조 필요
\end{itemize}

\note{\textbf{$k$를 로그 보간한 이유.} Si의 $k$는 250\,nm에서 3.67, 1100\,nm에서 $3\times10^{-5}$ — \textbf{5자릿수}에 걸칩니다. 간접갭이라 밴드 끝 근처 흡수가 포논 보조로만 일어나 지수적으로 죽기 때문입니다. 지수 꼬리를 선형 보간하면 격자점 사이에서 크게 틀립니다. 따라서 $n$은 선형, $\log k$는 선형 보간합니다.}

%==============================================================
\section{역문제 — 왜 어려운가}

\subsection{근본 축퇴}

투명막에서
\begin{equation}
R(d)\simeq A+B\cos(2\delta),\qquad \delta=\frac{2\pi}{\lambda}n_f d\cos\theta_f
\end{equation}
$R$은 $d$에 대해 주기적입니다. \textbf{서로 다른 두 축퇴}가 있으며, 해결책도 다릅니다.

\begin{center}
\begin{tabular}{@{}lll@{}}
\toprule
축퇴 & 원인 & 해결 \\
\midrule
Flank (측면) & $\cos$이 우함수 $\to$ 한 주기 안에서 같은 $R$이 두 번 & 위상 측정 (엘립소미트리) \\
Order (차수) & $\lambda/(2n_f\cos\theta_f)$마다 반복 & 광대역 또는 각도 스캔 \\
\bottomrule
\end{tabular}
\end{center}

실험 2의 설계 A(405 nm 단일점)에서 살아남은 후보 8개가 이 구조를 그대로 보여줍니다.
\begin{equation*}
\underbrace{57.1,\ 79.5}_{\text{flank 쌍}}\ \Big|\ \underbrace{194.9,\ 217.3}_{\text{flank 쌍}}\ \Big|\ \underbrace{332.8,\ 355.1}\ \Big|\ \underbrace{470.4,\ 492.9}
\end{equation*}
쌍 간 간격 137.82\,nm는 예측값 $\lambda/(2n_f)=137.79$\,nm와 0.02\% 일치합니다.

\note{\textbf{이것은 최적화 문제가 아닙니다.} 단일파장·단일각 반사율 측정은 원리적으로 $d$를 결정할 수 없고, 어떤 알고리즘·정칙화·평균화도 이를 고치지 못합니다. 정보가 데이터에 없습니다.}

\subsection{\texorpdfstring{$n\cdot d$ 축퇴 — 반사계의 원죄}{n·d 축퇴 — 반사계의 원죄}}

프린지 위상은 $n_f d$만 고정합니다. 따라서 가정한 굴절률이 틀리면
\begin{equation}
n_{\text{assumed}}\,d_{\text{fit}} = n_{\text{true}}\,d_{\text{true}}
\quad\Longrightarrow\quad
\boxed{\;\frac{\Delta d}{d} = +\frac{\Delta n}{n}\;},\qquad \Delta n \equiv n_{\text{true}}-n_{\text{assumed}}
\end{equation}

\note{\textbf{부호 방향에 주의.} 가정 굴절률이 \emph{낮으면} 적합 두께가 \emph{높게} 나옵니다. 실험 4의 S3에서 $\Delta n=+0.002$에 대해 측정 편향 $+0.3137$\,nm, 예측 $+0.2984$\,nm로 5\% 이내 일치를 확인했습니다. 한편 $d$와 $n$을 \emph{동시에} 적합하면 상관계수는 $\rho=-0.930$으로 음수입니다 — 축퇴 골짜기를 따라 움직일 때 $n$을 올리면 $d$를 내려야 $nd$가 유지되기 때문입니다. 두 부호는 서로 다른 상황을 기술하며 모순이 아닙니다.}

$\rho$가 정확히 $-1$이 아닌 유일한 이유는 \textbf{분산}입니다. $n_{\text{SiO}_2}(\lambda)$는 파장에 따라 변하는데 상수 오프셋은 변하지 않으므로 스펙트럼 형태에 미세한 차이가 남습니다. 협대역 측정에서는 이 차이마저 사라져 $n$과 $d$를 동시에 결정할 수 없습니다.

%==============================================================
\section{정밀도의 한계 — Fisher 정보량과 Cramér--Rao 하한}

\subsection{유도}

측정 $R_i$가 독립 가우시안 잡음 $\sigma$를 가지면
\begin{equation}
\ln L = -\frac12\sum_i\frac{(R_i^{\text{meas}}-R_i(d))^2}{\sigma^2}+\text{const}
\end{equation}
Fisher 정보량은
\begin{equation}
F_{ab} = \sum_i\frac{1}{\sigma_i^2}\frac{\partial R_i}{\partial p_a}\frac{\partial R_i}{\partial p_b},
\qquad \operatorname{Cov}(\hat p)\succeq F^{-1}
\end{equation}
단일 파라미터에서는
\begin{equation}
\boxed{\;\sigma_d \ \ge\ \frac{\sigma_R}{\sqrt{\sum_i(\partial R_i/\partial d)^2}}\;}
\end{equation}

\subsection{물리적 읽기}

\begin{equation}
\frac{\partial R}{\partial d} = -2B\cdot\frac{2\pi n_f\cos\theta_f}{\lambda}\sin2\delta
\end{equation}

세 가지가 즉시 보입니다.

\begin{enumerate}[leftmargin=16pt,itemsep=1pt]
\item \textbf{감도 $\propto 1/\lambda$} — 405\,nm가 633\,nm보다 유리합니다 (실험 1: $6.90$ vs $3.90\times10^{-3}$/nm).
\item \textbf{감도 $\propto B$ (프린지 대비)} — SiO$_2$/Si는 $r_{01}\approx-0.19$, $r_{12}\approx-0.45$로 크기가 비슷해 대비가 좋습니다.
\item \textbf{$\sin2\delta=0$인 두께에서 감도 0} — 반주기(405 nm에서 68.9\,nm)마다 \textbf{사각지대}가 존재합니다.
\end{enumerate}

\subsection{가장 중요한 비대칭}

\note{\textbf{CRLB는 분산(랜덤 오차)만 제한합니다. 편향(계통 오차)에 대해서는 아무 말도 하지 않습니다.} 적합 결과가 CRLB에 정확히 앉아 있으면서 동시에 1\,nm 틀릴 수 있습니다. 실험 3(랜덤)과 실험 4(계통)를 끝까지 분리해 서술하는 이유가 이것입니다. 실측된 효율은 $\text{CRLB}/\sigma_{\text{MC}} = 1.019$로, 추정기가 데이터의 정보를 사실상 전부 뽑아냅니다 — 즉 \textbf{더 좋아지려면 알고리즘이 아니라 측정을 바꿔야 합니다}.}

\subsection{계통오차의 물리}

\textbf{계면층 (S1).} 열산화 SiO$_2$는 결정 Si와 급격히 만나지 않고, 아준화학량론 SiO$_x$ 전이 영역이 존재합니다. 얇은 층은 $r_{12}$의 \emph{크기와 위상을 모두} 섭동하는데, 피팅 모델은 $d$ 하나만 조절할 수 있으므로 위상 성분을 $d$로 흡수합니다. 전달계수는 0.745\,nm/nm입니다. \textbf{가장 위험한 점은 잔차 RMS가 $3\times10^{-4}$에 머문다는 것입니다} — 좋은 $\chi^2$가 답을 보증하지 않습니다.

\textbf{거칠기 (S2).} 특징 크기 $\ll\lambda$이면 준정적 근사가 성립해 유효 매질로 대체할 수 있습니다. Bruggeman 자기무결 조건은
\begin{equation}
f_a\frac{\varepsilon_a-\varepsilon}{\varepsilon_a+2\varepsilon}+f_b\frac{\varepsilon_b-\varepsilon}{\varepsilon_b+2\varepsilon}=0
\;\Longrightarrow\;
2\varepsilon^2-b\varepsilon-\varepsilon_a\varepsilon_b=0,\quad
b=(3f_a-1)\varepsilon_a+(2-3f_a)\varepsilon_b
\end{equation}

\note{\textbf{왜 Maxwell--Garnett이 아닌가.} MG는 호스트/내포물 구분을 요구합니다. 50:50에서는 둘 중 무엇이 호스트인지 정할 수 없으므로 자기무결(self-consistent) Bruggeman이 맞습니다.\\[4pt]
\textbf{그리고 이건 측정 오차가 아니라 정의 문제입니다.} 측정된 전달계수 0.438\,nm/nm는 피팅이 물리적 표면 위치가 아니라 \textbf{부피등가 두께} $d+f\,t_{\text{rough}}$를 보고한다는 뜻입니다. 거칠기 규약이 다른 두 장비는 둘 다 완벽히 작동해도 $\sim f\,t_{\text{rough}}$만큼 불일치합니다. 실무에서 장비 간 매칭이 안 되는 흔한 원인입니다.}

\textbf{참조 데이터 (S4).} Green 2008과 Aspnes 1983은 둘 다 동료심사를 거쳤고 어느 쪽도 틀리지 않았습니다. 405\,nm에서 $\Delta n=+0.034$, 633\,nm에서 $-0.009$입니다. 제거 불가능하며 \textbf{계통 불확도로 안고 가야} 합니다.

\subsection{최종 예산}

\begin{center}
\begin{tabular}{@{}lrr@{}}
\toprule
원인 & 편향 (nm) & 랜덤 바닥 대비 \\
\midrule
S1 계면 SiO$_x$ 1 nm & $+0.745$ & 168 \\
S2 거칠기 2 nm (Bruggeman) & $+0.876$ & 197 \\
S3 막 굴절률 $\Delta n=+0.004$ & $+0.627$ & 141 \\
S4 Si 참조데이터 (Green $\leftrightarrow$ Aspnes) & $+0.166$ & 37 \\
S5 각도 오차 $0.1°$ & $-0.147$ & 33 \\
\midrule
\textbf{RSS 합성} & \textbf{1.33} & \textbf{299} \\
\bottomrule
\end{tabular}
\end{center}

\begin{equation*}
\underbrace{4\ \text{pm}}_{\text{알고리즘 바닥}}\quad\text{vs}\quad\underbrace{1.33\ \text{nm}}_{\text{현실 정확도}}\qquad\text{비율}\approx300
\end{equation*}

개선 방향은 옵티마이저가 아니라 \textbf{시료 지식}입니다.

\newpage

\part*{\color{accent}Part II — 주석 달린 소스}
\addcontentsline{toc}{section}{Part II — 주석 달린 소스}

\section{\texttt{omf/core.py} — 순방향 모델}

이 모듈이 프로젝트의 심장입니다. 259줄 중 실제 계산은 약 60줄이고, 나머지는 규약과
근거의 문서화입니다. 물리적으로 되돌리기 어려운 결정(규약, 분지, 플럭스 정의)이
전부 여기에 모여 있으므로 그렇게 설계했습니다.

\subsection{규약 선언과 분지 상수}
\lstinputlisting[firstnumber=1]{frag/c001.py}
\note{
\textbf{왜 규약을 파일 최상단에 못 박는가.} $e^{-i\omega t}$, $\nt=n+ik$, $e^{+ik_z z}$
세 선택은 독립이 아니라 하나가 정해지면 나머지가 강제됩니다(§3.1). 문헌의 절반은
$e^{+i\omega t}$, $\nt=n-ik$를 쓰며 둘 다 옳지만, 한 계산 안에서 섞으면 반드시 틀립니다.
이 파일에서 규약은 여기서 한 번 선언되고 이후 어디서도 재론되지 않습니다.\\[4pt]
\texttt{\_BRANCH\_TOL = 1e-10}은 ``무손실 전파''를 판정하는 수치 임계입니다.
$\Imag(k_z)$가 정확히 0인 경우는 부동소수점에서 거의 나오지 않으므로, 이 폭 안에
들어오면 무손실로 간주하고 대신 $\Real(k_z)>0$을 요구합니다.
}
\subsection{분지 선택 — 이 코드의 유일한 물리적 강제}
\lstinputlisting[firstnumber=47]{frag/c002.py}
\note{
\textbf{이 8줄이 특수 케이스를 전부 없앱니다.} $+z$ 진행파는 $e^{i\Real(k_z)z}e^{-\Imag(k_z)z}$
이므로 수동 매질에서 감쇠하려면 $\Imag(k_z)\ge0$이어야 합니다. 반대 분지는
(i) 거리에 따라 \emph{증폭}되는 비물리적 파를 기술하고,
(ii) 두꺼운 흡수층에서 $e^{ik_zd}$가 오버플로를 일으켜 계산을 파괴합니다.\\[4pt]
\textbf{두 번째 조건 \texttt{(|Im| <= tol) \& (Re < 0)}의 역할.} 무손실 매질에서는
$\Imag(k_z)=0$이라 첫 조건이 아무것도 걸러내지 못합니다. 이때는 부호를 정할 다른 기준이
필요한데, 그것이 ``순방향 = $+z$''라는 정의, 즉 $\Real(k_z)>0$입니다. 이 줄이 없으면
무손실 층에서 전·후향파가 뒤바뀌는 미묘한 버그가 생깁니다.\\[4pt]
\texttt{np.where}로 쓴 이유는 \texttt{if}가 아니라 \emph{배열 전체}에 대해 동시에
적용해야 하기 때문입니다. 이것이 뒤의 벡터화를 가능하게 합니다.
}
\subsection{면내 파수 보존 = 스넬 법칙}
\lstinputlisting[firstnumber=66]{frag/c003.py}
\note{
\textbf{$k_x$를 보존량으로 쓴 것이 이 코드의 가장 중요한 설계 결정입니다.} 면내 병진
대칭성으로부터 $k_x = n_0k_0\sin\theta_0$이 모든 층에서 같고,
$k_{z,j}=\sqrt{(\nt_jk_0)^2-k_x^2}$입니다. 이것이 스넬 법칙 \emph{그 자체}입니다.\\[4pt]
각도 $\theta_j=\arcsin(k_x/\nt_jk_0)$로 쓰면 흡수 매질에서 $\theta_j$가 복소수가 되고,
복소 arcsin이 어느 리만 시트에 있는지를 매번 추적해야 합니다. $k_z$로 쓰면
전반사·흡수·소멸파가 전부 ``$k_z$가 복소수가 된다''는 한 문장으로 흡수됩니다.
\textbf{if 분기가 0개인 코드}가 나온 이유입니다.\\[4pt]
\texttt{snell\_angles}가 존재하지만 솔버는 이를 한 번도 호출하지 않습니다. 보고용입니다 —
사람이 읽을 때는 각도가 직관적이기 때문입니다.\\[4pt]
\textbf{브로드캐스팅 세부.} \texttt{n} 은 \texttt{(..., N)}, \texttt{k0}와 \texttt{kx}는
\texttt{(...)} 모양이므로 \texttt{[..., None]}로 축을 하나 늘려 층 축과 정렬시킵니다.
이 shape 규약이 모듈 전체를 관통합니다.
}
\subsection{Fresnel 계수}
\lstinputlisting[firstnumber=102]{frag/c004.py}
\note{
\textbf{$\cos\theta$ 대신 $k_z$로 쓴 효과.} $\nt_m\cos\theta_m = k_{zm}/k_0$이므로 모든 비율에서
공통 인자 $k_0$가 소거됩니다. 결과적으로 s편광은
$r=(k_{zi}-k_{zj})/(k_{zi}+k_{zj})$라는 임피던스 정합 형태가 되고, 이는 전송선의
$\Gamma=(Z_2-Z_1)/(Z_2+Z_1)$과 글자 그대로 같습니다(§2.2).\\[4pt]
p편광의 $\nt^2$ 인자는 $\mathbf{E}$가 아니라 $\mathbf{H}$가 접선 성분이기 때문에 나옵니다.
$E_x\propto(k_z/\varepsilon)H_y$이므로 어드미턴스가 $Y_p\propto\nt^2/k_z$로 뒤집힙니다.
$t_p$의 $2\nt_i\nt_j$는 $H$ 진폭 결과를 $E$ 진폭으로 환산한 것입니다.\\[4pt]
\textbf{$r_p$의 부호는 규약입니다.} 반사 후 $\mathbf{E}_p$의 양의 방향을 어떻게 정하느냐에
따라 달라지며, 이 코드는 Byrnes 규약($\theta_B$ 아래에서 $r_p>0$)을 씁니다.
$|r_p|^2$는 규약 무관이므로 $R_p$는 안전합니다. 다른 출처의 $r_p$를 섞어 쓰면 사고가 납니다.\\[4pt]
인자로 배열 슬라이스 \texttt{n[..., :-1]}와 \texttt{n[..., 1:]}를 받도록 설계해,
\textbf{모든 계면쌍을 호출 한 번에} 처리합니다.
}
\subsection{전달행렬 조립 — 함수의 골격}
\lstinputlisting[firstnumber=140]{frag/c005.py}
\note{
\textbf{입력 검증이 물리 조건입니다.} 바깥 두 층의 두께가 \texttt{inf}여야 하는 것은
반무한 매질 가정이고, 입사 매질이 실수 굴절률이어야 하는 것은 $R$과 $T$가
플럭스 비로 정의되려면 입사파 자체가 감쇠하지 않아야 하기 때문입니다.
흡수하는 매질에서 ``입사 플럭스''는 위치에 따라 달라져 정의되지 않습니다.\\[4pt]
\texttt{np.broadcast\_shapes(th\_0.shape, lam\_vac.shape, n\_list.shape[:-1])}가 이 모듈의
벡터화 핵심입니다. 파장 배열 $(M,)$과 각도 배열 $(K,)$를 각각
$(M,1)$, $(1,K)$로 넣으면 batch가 $(M,K)$가 되어 \textbf{2차원 스캔이 호출 한 번}에 끝납니다.
실험 1의 각도 스캔과 실험 2의 설계 B가 이 성질을 씁니다.
}
\subsection{위상과 계면 항 — 벡터화가 일어나는 지점}
\lstinputlisting[firstnumber=187]{frag/c006.py}
\note{
\textbf{여기서 슬라이싱 두 줄이 전체 스택을 처리합니다.}
\texttt{delta}는 내부층만 $(\ldots,N-2)$, \texttt{r\_if/t\_if}는 계면 수만큼 $(\ldots,N-1)$입니다.
층 인덱스가 배열 축으로 올라가 있기 때문에 파이썬 루프가 필요 없습니다.\\[4pt]
\textbf{$\Imag(\delta)\ge0$이 보장되는 것이 안전성의 근거입니다.} 분지 규칙 덕분에
$e^{+i\delta}$의 크기가 1을 넘지 않으므로, 두께가 아무리 커도 오버플로가 발생하지 않습니다.
분지 교정을 생략하면 정확히 여기서 계산이 터집니다.
}
\subsection{층 행렬과 누적 곱}
\lstinputlisting[firstnumber=202]{frag/c007.py}
\note{
\textbf{루프가 데이터가 아니라 층을 돕니다.} 이것이 196배 가속의 정체입니다.
층은 보통 3$\sim$5개인데 데이터 점은 수천 개이므로, 루프를 어디에 두느냐가 성능을 결정합니다.
\texttt{einsum("...ij,...jk->...ik")}가 batch 차원 전체에 대해 $2\times2$ 곱을 한 번에 수행합니다.\\[4pt]
\textbf{$D$와 $P$의 물리적 역할이 분리되어 있습니다.} $D$(계면)는 전·후향 진폭을 섞고,
$P$(전파)는 위상만 축적합니다. Airy 급수를 무한히 더하는 대신 이 두 연산을 번갈아
곱하는 것이 전달행렬의 전부입니다 — 급수를 미리 합해둔 형태입니다(§4.3).\\[4pt]
\textbf{$r=M_{21}/M_{11}$, $t=1/M_{11}$의 근거.} 경계조건
$[E_0^+,E_0^-]^T = M[E_f^+,0]^T$에서 출사 매질에 후향파가 없다($E_f^-=0$)는 조건을 쓰면
$E_0^+=M_{11}E_f^+$, $E_0^-=M_{21}E_f^+$이므로 두 식이 즉시 따라옵니다.
$N=3$에 대해 직접 곱해보면 정확히 Airy 공식이 나오며, 검증 B5가 이를 $2\times10^{-16}$로 확인합니다.
}
\subsection{반사율과 투과율 — p편광의 켤레}
\lstinputlisting[firstnumber=237]{frag/c008.py}
\note{
\textbf{$T$는 $|t|^2$가 아닙니다.} $z$방향 시간평균 Poynting 플럭스의 비입니다.
s편광은 $S_z\propto\Real(k_z)|E|^2$이라 인자가 $\Real(k_{zf})/\Real(k_{z0})$이지만,
p편광은 $S_z\propto\Real(\nt^2k_z^*)|E|^2/|\nt|^2$이라 \textbf{켤레의 위치가 다릅니다}.\\[4pt]
$\nt$가 실수면 두 식이 일치하므로 무손실 문제만 다루면 이 차이를 영영 발견하지 못합니다.
흡수 매질에서 p에 s식을 쓰면 에너지 보존이 깨지고, 오차는 각도와 $\Imag(\nt)$에 따라 커집니다.
검증 A3가 이를 잡아냅니다.\\[4pt]
\textbf{$A = 1-R-T$의 의미에 주의.} 반무한 흡수 기판에서는 $T$가 ``기판 \emph{안으로 들어간}''
플럭스이므로 $R+T=1$이고 $A=0$이 됩니다. 기판 내부의 흡수는 $A$가 아니라 $T$에 계산됩니다.
불투명성은 \textbf{유한 슬랩}의 성질입니다(§5.1, 검증 B12/B13).
}

\section{\texttt{omf/materials.py} — 굴절률 공급}

두 종류의 모델이 공존하며, 그 이유가 물리적으로 명확합니다. 투명 영역에서는 해석식
(Sellmeier)이 정확하고, 강한 밴드간 흡수가 있으면 해석식이 존재하지 않아 측정 표를
보간해야 합니다.

\subsection{Sellmeier — 감쇠 없는 Lorentz 진동자}
\lstinputlisting[firstnumber=34]{frag/c009.py}
\note{
\textbf{Sellmeier는 피팅 공식이 아니라 물리 모델입니다.} Lorentz 진동자 합
$\varepsilon=1+\sum_j f_j\omega_{pj}^2/(\omega_j^2-\omega^2-i\gamma_j\omega)$에서
투명 영역($\gamma\to0$) 극한을 취하고 $\omega\to\lambda$로 바꾼 것이 정확히
$n^2=1+\sum_j B_j\lambda^2/(\lambda^2-C_j)$입니다.\\[4pt]
Malitson 계수의 $\sqrt{C_j}$는 각각 68.4\,nm(18.1\,eV, 심준위), 116.2\,nm(10.7\,eV,
SiO$_2$ 밴드갭/엑시톤), 9.896\,$\mu$m(125\,meV, Si--O--Si 신축 포논)입니다.
\textbf{가시광 SiO$_2$의 분산은 UV 전자 전이와 IR 포논 사이에 끼인 꼬리}입니다.\\[4pt]
$k=0$은 구성상 그렇게 되는 것이지 근사가 아닙니다 — $\gamma\to0$을 이미 가정했기 때문입니다.
융용실리카의 가시광 흡수는 $\alpha<10^{-4}$/cm이라 이 가정이 훌륭합니다.\\[4pt]
\textbf{docstring에 적힌 경고가 실험 4의 S3로 이어집니다.} Si 위에 성장한 \emph{열산화}
SiO$_2$는 벌크 융용실리카와 같지 않습니다. 치밀화와 응력으로 가시광 $n$이 보통
0.002$\sim$0.006 높습니다. 이를 무시하면 217\,nm 막에서 0.6\,nm 편향이 생깁니다.
}
\subsection{비분산 매질과 상수 공장}
\lstinputlisting[firstnumber=61]{frag/c010.py}
\note{
\texttt{n\_air}가 분산을 무시하는 근거는 정량적입니다 — 가시광 전역에서 공기의 분산은
$\sim3\times10^{-6}$으로, 이 프로젝트가 분해할 수 있는 어떤 효과보다도 작습니다.
근사를 쓸 때는 ``작다''가 아니라 \textbf{무엇보다 작은지}를 밝혀야 합니다.\\[4pt]
\texttt{n\_bk7}의 docstring은 실제 커버슬립이 D263이나 B270이라 $n_d$가 1.5230 대 1.5168로
약 0.4\% 다르다고 명시합니다. Brewster 각 시연에는 충분하지만 절대 두께 계측에는 부족합니다 —
이런 구분을 코드 안에 남겨두는 것이 나중에 자신을 속이지 않는 방법입니다.\\[4pt]
\texttt{n\_constant}는 클로저를 반환하는 공장 함수입니다. 실험 4에서 계면층
$n=2.5$를 끼워 넣을 때 쓰이며, 덕분에 ``물질''을 이름·함수·상수 어느 것으로도
지정할 수 있는 통일 인터페이스가 됩니다.
}
\subsection{Si — 표 보간과 로그 스케일}
\lstinputlisting[firstnumber=94]{frag/c011.py}
\note{
\textbf{왜 Si에는 해석식이 없는가.} 가시·UV에서 밴드간 흡수가 강해 $\gamma\to0$ 가정이
붕괴합니다. 임계점은 $E_1\approx3.4$\,eV(365\,nm, $\Lambda$ 방향)와
$E_2\approx4.25$\,eV(292\,nm, $X/\Sigma$ 근처)이며, $E_1$에서 $n$이 6.03까지 치솟습니다.
저차 다항식이나 Sellmeier로는 이 구조를 재현할 수 없습니다.\\[4pt]
\textbf{$\log k$ 보간이 핵심입니다.} $k$는 250\,nm의 3.67에서 1100\,nm의 $3\times10^{-5}$까지
\emph{5자릿수}에 걸칩니다. Si의 갭이 간접갭이라 밴드 끝 근처 흡수는 포논 보조로만 일어나
지수적으로 죽기 때문입니다. 지수 꼬리를 선형 보간하면 격자점 사이에서 크게 틀립니다.
$n$은 완만하므로 선형으로 충분합니다.\\[4pt]
\textbf{범위 밖 예외를 던지는 이유.} \texttt{np.interp}는 범위 밖에서 조용히 끝점 값을
반환합니다. 계측 코드에서 이는 최악의 실패 모드입니다 — 틀린 답이 경고 없이 나옵니다.
명시적으로 예외를 던져 사용자가 알아채도록 강제합니다.
}
\subsection{두 번째 Si 데이터셋 — 계통오차를 재기 위한 장치}
\lstinputlisting[firstnumber=142]{frag/c012.py}
\note{
\textbf{Aspnes 데이터는 편의 기능이 아니라 실험 장비입니다.} Green 2008과
Aspnes 1983은 둘 다 동료심사를 거쳤고 \emph{어느 쪽도 틀리지 않았습니다}.
405\,nm에서 $\Delta n=+0.034$, 633\,nm에서 $-0.009$ 차이가 납니다.\\[4pt]
이 차이는 수치 잡음이 아니라 \textbf{실재하는 참조 데이터 불일치}이며, 독립 측정 없이는
제거할 수 없습니다. 두 데이터셋을 모두 탑재해 한쪽으로 생성하고 다른 쪽으로 적합하면
그 불일치가 두께 편향 몇 nm에 해당하는지 정량화됩니다(실험 4의 S4: 0.166\,nm).
\textbf{제거하는 대신 예산에 실어 나르는 것}이 정직한 처리입니다.\\[4pt]
\texttt{build\_n\_list}는 이름·콜러블·복소상수를 모두 받아 \texttt{(..., N)} 배열로
쌓습니다. 이 통일 인터페이스 덕분에 실험 4에서 층 하나를 끼워 넣거나 굴절률에
오프셋을 더하는 조작이 한 줄로 끝납니다.
}

\section{\texttt{omf/fitter.py} — 역문제}

순방향 모델이 정확해도 역문제는 별개의 어려움을 갖습니다. 이 모듈의 존재 이유는
\textbf{$R(d)$가 $d$에 대해 주기적}이라는 사실 하나입니다.

\subsection{모듈 서두 — 왜 단순 최소제곱이 실패하는가}
\lstinputlisting[firstnumber=1]{frag/c013.py}
\note{
\textbf{프린지 차수 축퇴는 알고리즘 결함이 아니라 정보의 부재입니다.}
$R\simeq A+B\cos2\delta$이므로 $\Delta d=\lambda/(2n_f\cos\theta_f)$만큼 두꺼운 막은
단일 파장에서 거의 같은 반사율을 냅니다(405\,nm에서 137.79\,nm).
비용함수가 빗살 모양이 되고, 경사법은 시작점에서 가장 가까운 빗살 이빨로 떨어집니다.\\[4pt]
서두에서 두 해결책을 명확히 구분합니다. (1) 조밀 격자 사전탐색은 \emph{옵티마이저}
문제를 풀고, (2) 광대역/각도 데이터는 \emph{정보} 문제를 풉니다. 후자가 없으면
전자로도 못 고칩니다 — 실험 2의 설계 A가 이를 실증합니다.
}
\subsection{FitResult — 불확도를 결과에 동봉}
\lstinputlisting[firstnumber=46]{frag/c014.py}
\note{
\textbf{추정값과 불확도를 한 객체로 묶은 것이 설계 의도입니다.} 계측에서 불확도 없는
숫자는 정보가 아닙니다. \texttt{correlation} 속성이 별도로 있는 이유는
$d$--$n$ 상관계수가 이 문제의 핵심 진단량이기 때문입니다(실험 3에서 $\rho=-0.930$).\\[4pt]
상관행렬을 계산할 때 \texttt{errstate}로 0-나눗셈을 막고 \texttt{nan}으로 두는 것은,
특이한 공분산(정보가 없는 방향)을 조용히 0으로 만들지 않기 위해서입니다.
}
\subsection{FilmModel 생성자 — 무엇을 자유 파라미터로 둘 것인가}
\lstinputlisting[firstnumber=74]{frag/c015.py}
\note{
\texttt{fit\_dn} 옵션이 이 클래스의 물리적 긴장 지점입니다. 열산화막과 벌크 융용실리카의
굴절률 차이를 흡수하려면 $n$을 자유롭게 두어야 하지만, \textbf{프린지가 실제로 재는 것은
광학 두께 $n\cdot d$}이므로 $d$와 $n$이 강하게 반상관됩니다.\\[4pt]
docstring이 이 함정을 미리 경고합니다. 실측된 결과는 $\rho=-0.930$이고
$\sigma_d$가 2.7배 악화됩니다. $\rho$가 정확히 $-1$이 아닌 유일한 이유는 분산입니다 —
$n_{\text{SiO}_2}(\lambda)$는 파장에 따라 변하는데 상수 오프셋은 변하지 않으므로
스펙트럼 형태에 미세한 차이가 남습니다. \textbf{협대역 측정에서는 이 차이마저 사라져
$n$과 $d$를 동시에 결정할 수 없습니다.}\\[4pt]
생성자 검증(바깥 층 \texttt{inf}, 적합 대상 층은 유한)은 물리적으로 말이 되는 스택만
허용하기 위한 것입니다.
}
\subsection{순방향 래퍼}
\lstinputlisting[firstnumber=115]{frag/c016.py}
\note{
\texttt{'unpol'}은 $R_s$와 $R_p$의 \textbf{강도 평균}입니다. 진폭 평균이 아닙니다 —
두 편광은 서로 코히런트하지 않으므로 간섭하지 않고, 검출기가 보는 것은 전력의 합입니다.
이를 진폭에서 평균하면 존재하지 않는 간섭항이 생깁니다.\\[4pt]
\texttt{n\_list}를 매 호출마다 \texttt{build\_n\_list}로 재생성하는 것은 비효율처럼 보이지만,
분산 모델이 파장에만 의존하므로 캐싱하면 \texttt{fit\_dn} 경로에서 오염이 생깁니다.
\texttt{.copy()} 후 오프셋을 더하는 것도 같은 이유입니다.
}
\subsection{조밀 격자 사전탐색 — 전역성 확보}
\lstinputlisting[firstnumber=141]{frag/c017.py}
\note{
\textbf{이 무식한 전수 탐색이 가능한 이유가 벡터화입니다.} 격자점 하나마다
\texttt{reflectance}를 호출하지만, 그 호출 자체가 이미 모든 파장에 대해 벡터화되어 있습니다.
1200점 격자 $\times$ 251 파장이면 30만 회 TMM 평가인데 1초 미만입니다.
\texttt{core.py}에서 루프를 층으로 옮긴 결정이 여기서 배당을 지급합니다.\\[4pt]
\textbf{실측된 효과.} 소박한 국소 최적화는 9개 시작점 중 \textbf{6개}에서 틀린 프린지
차수로 수렴했습니다(47.0, 388.2, 506.2\,nm 등). 사전탐색을 붙이면 9개 전부 성공합니다.\\[4pt]
격자 간격은 프린지 주기보다 충분히 촘촘해야 합니다 — 137.79\,nm 주기에 대해
1200점/600\,nm는 0.5\,nm 간격이므로 안전합니다.
}
\subsection{정밀화와 공분산}
\lstinputlisting[firstnumber=166]{frag/c018.py}
\note{
\textbf{공분산 추정의 두 갈래가 통계적으로 다른 주장입니다.}\\[3pt]
$\bullet$ \texttt{sigma}를 주면 잔차가 이미 백색화되어 있으므로 $(J^TJ)^{-1}$이 곧
공분산입니다. 이는 \emph{사전} 추정 — ``내 장비의 잡음이 이만큼임을 독립적으로 안다''.\\[3pt]
$\bullet$ \texttt{sigma=None}이면 축소 카이제곱 $\chi^2_{\text{red}}=\text{SSR}/(m-p)$을
곱합니다. 이는 \emph{사후} 추정 — ``잔차 크기로부터 잡음을 역산한다''. 모델이 틀렸을 때
$\chi^2_{\text{red}}$가 커져 불확도가 부풀려지므로, 어느 정도 자기 보호 기능이 있습니다.\\[4pt]
\textbf{그러나 이 보호는 불완전합니다.} 실험 4의 S1에서 1\,nm 계면층은 0.745\,nm 편향을
만들지만 잔차 RMS는 $3\times10^{-4}$에 머뭅니다. \textbf{좋은 $\chi^2$가 답을 보증하지 않습니다.}
계통오차는 잔차에 거의 흔적을 남기지 않고 파라미터로 흡수되기 때문입니다.\\[4pt]
\texttt{xtol/ftol/gtol}을 $10^{-14}$로 조인 것은 실험 3에서 CRLB와 비교하기 위함입니다.
수렴 허용오차가 헐거우면 몬테카를로 산포에 옵티마이저 잡음이 섞여 효율이 1보다
작게 나옵니다 — 측정하려는 양이 오염됩니다.
}

\section{\texttt{omf/uncertainty.py} — 판정}

이 모듈은 새로운 숫자를 만들지 않습니다. 앞의 모듈들이 만든 숫자가 \textbf{믿을 만한지}를
판정합니다. 계측 프로젝트와 단순 시뮬레이션을 가르는 지점입니다.

\subsection{모듈 서두 — 왜 CRLB인가}
\lstinputlisting[firstnumber=1]{frag/c019.py}
\note{
\textbf{두 질문이 다릅니다.} (1) 이 측정 설계에 정보가 얼마나 들어 있는가(CRLB) —
\emph{실험}의 성질입니다. (2) 내 추정기가 그것을 다 뽑아내는가(몬테카를로) —
\emph{코드}의 성질입니다. 둘을 비교해야 어디를 고쳐야 하는지 알 수 있습니다.\\[4pt]
서두 마지막 문단이 이 프로젝트에서 가장 중요한 문장입니다.
\textbf{CRLB는 분산만 제한하고 편향에 대해서는 아무 말도 하지 않습니다.}
적합 결과가 CRLB에 정확히 앉아 있으면서 동시에 3\,nm 틀릴 수 있습니다 —
시료에 모델이 빠뜨린 자연 산화막이 있으면 그렇습니다.
랜덤 정밀도와 계통 정확도는 별개의 예산입니다.
}
\subsection{수치 야코비안과 Fisher 정보량}
\lstinputlisting[firstnumber=40]{frag/c020.py}
\note{
\textbf{중심차분을 쓴 이유.} 절단오차가 $O(h^2)$이라 전방차분의 $O(h)$보다 한 자릿수
유리합니다. $h$를 파라미터 스케일의 $\sim10^{-4}$로 잡으면 절단오차가 고려하는
잡음 수준보다 충분히 아래로 내려갑니다. 두께에 대해 \texttt{steps=[0.02]}\,nm를 쓴 것이
이 판단입니다 — 프린지 주기 137.79\,nm의 $1.5\times10^{-4}$입니다.\\[4pt]
\textbf{Fisher 행렬의 물리적 읽기.} $F=J^TWJ$에서
$\partial R/\partial d = -2B\cdot(2\pi n_f\cos\theta_f/\lambda)\sin2\delta$이므로 세 가지가 보입니다:
감도 $\propto1/\lambda$(405\,nm가 633\,nm보다 유리),
감도 $\propto B$(프린지 대비),
그리고 $\sin2\delta=0$인 두께마다 \textbf{감도가 정확히 0인 사각지대}가 반주기 간격으로 존재합니다.\\[4pt]
$F$의 역행렬 대각선이 각 파라미터의 하한 분산입니다. 다변수에서 $(F^{-1})_{dd}$가
$1/F_{dd}$보다 \emph{큰} 것이 파라미터 상관의 대가입니다.
}
\subsection{몬테카를로 — 그리고 inverse crime의 명시}
\lstinputlisting[firstnumber=73]{frag/c021.py}
\note{
\textbf{판정 기준.} 산포가 CRLB와 같으면 추정기가 효율적이고, 크면 정보를 흘리고 있으며,
평균이 참값과 다르면 편향이 있습니다. 실측 결과는 세 자릿수 잡음 범위에서
평균 효율 \textbf{1.019}, 편향 $<0.005$\,nm였습니다.\\[4pt]
\textbf{docstring의 마지막 문단이 이 프로젝트의 정직성 장치입니다.} 같은 모델이 데이터를
생성하고 적합하므로(inverse crime), 존재하는 오차원은 주입한 잡음과 옵티마이저뿐입니다.
틀린 분산 모델, 무시된 자연 산화막, 거칠기, 유한 빔 발산 — \textbf{모든 모델 형태 오차가
구성상 보이지 않습니다}.\\[4pt]
따라서 여기서 나온 $\sigma$는 \emph{잡음 바닥}, 즉 성능의 상한이지 검증된 정확도가
아닙니다. 이 경고를 코드 안에 박아 둔 이유는, 나중에 결과만 인용할 때 스스로를
속이지 않기 위해서입니다. 실험 4가 이 경고를 정량화합니다 — 4\,pm 대 1.3\,nm, 300배.
}

\section{실험 스크립트가 조작하는 손잡이}

네 스크립트는 같은 코어를 쓰되 \textbf{서로 다른 변수를 흔듭니다}. 무엇을 흔드는지가
각 실험이 답하는 질문을 정의합니다.

\begin{center}
\begin{tabular}{@{}llp{7.6cm}@{}}
\toprule
파일 & 흔드는 손잡이 & 답하는 질문 \\
\midrule
exp01 & $d$, $\lambda$, $\theta$ & 반사율이 어떻게 생겼고, 어디에 두께 정보가 있는가 \\
exp02 & \textbf{측정 설계} & 이 데이터로 $d$를 유일하게 정할 수 있는가 \\
exp03 & \textbf{잡음 $\sigma_R$} & 정보 한계에 도달했는가 \\
exp04 & \textbf{생성 모델 $\neq$ 적합 모델} & 모델이 틀리면 얼마나 틀리는가 \\
\bottomrule
\end{tabular}
\end{center}

exp01$\sim$03은 같은 모델로 데이터를 만들고 맞춥니다(inverse crime).
\textbf{exp04만 이를 깹니다} — 생성 쪽에 계면층이나 거칠기층을 끼우거나, 다른 Si
데이터셋을 쓰거나, 각도를 틀리게 줍니다. 그 세 줄이 4\,pm과 1.3\,nm를 갈라놓습니다.

\subsection{inverse crime을 깨는 코드}
\begin{lstlisting}[firstnumber=1]
# exp04: 생성은 4층, 적합은 3층
n_true = build_n_list(["air", "SiO2", n_constant(2.5), "Si"], lamC)
R      = coh_tmm("s", n_true, [INF, D_TRUE, t, INF], thC, lamC)["R"]
fr     = fit_to(R)          # fit_model 은 air/SiO2/Si 3층 고정
bias   = fr.params[0] - D_TRUE
\end{lstlisting}
\note{
\textbf{4줄이 프로젝트의 결론을 만듭니다.} 생성 스택에 층 하나를 끼워 넣고 적합 모델은
그대로 두면, 적합기는 계면층이 만든 위상 섭동을 \emph{조절 가능한 유일한 파라미터}인
$d$로 흡수할 수밖에 없습니다. 1\,nm 계면층 $\to$ 0.745\,nm 편향, 랜덤 바닥의 168배입니다.\\[4pt]
같은 패턴으로 거칠기(Bruggeman EMA 층 추가), 굴절률 오프셋(생성 시 $+\Delta n$),
참조 데이터(생성은 Aspnes, 적합은 Green), 각도 오차(생성 각도에 $+\delta$)를 각각
격리해 측정합니다. \textbf{한 번에 하나씩 흔드는 것}이 요점입니다 — 그래야 예산의
각 항목이 무엇에서 왔는지 추적됩니다.
}

\part*{\color{accent}Part III — 전체 원본 소스}
\addcontentsline{toc}{section}{Part III — 전체 원본 소스}

\section{축약 없는 소스 코드}

Part II에서 지면상 접었던 docstring을 포함해 저장소의 모든 소스를 그대로 싣습니다.
줄 번호는 각 파일의 실제 줄 번호입니다.
\subsection{\texttt{omf/core.py} — TMM 엔진 (260줄)}
\lstinputlisting[firstnumber=1]{/home/claude/optical-metrology-fitter/omf/core.py}
\subsection{\texttt{omf/materials.py} — 분산 모델 (191줄)}
\lstinputlisting[firstnumber=1]{/home/claude/optical-metrology-fitter/omf/materials.py}
\subsection{\texttt{omf/fitter.py} — 역문제 솔버 (222줄)}
\lstinputlisting[firstnumber=1]{/home/claude/optical-metrology-fitter/omf/fitter.py}
\subsection{\texttt{omf/uncertainty.py} — 불확도 분석 (105줄)}
\lstinputlisting[firstnumber=1]{/home/claude/optical-metrology-fitter/omf/uncertainty.py}
\subsection{\texttt{omf/\_\_init\_\_.py} — 패키지 인터페이스 (22줄)}
\lstinputlisting[firstnumber=1]{/home/claude/optical-metrology-fitter/omf/__init__.py}
\subsection{\texttt{tests/test\_core.py} — 검증 스위트 (34 항목) (416줄)}
\lstinputlisting[firstnumber=1]{/home/claude/optical-metrology-fitter/tests/test_core.py}
\subsection{\texttt{scripts/exp01\_forward\_model.py} — 실험 1 — 순방향 모델 (181줄)}
\lstinputlisting[firstnumber=1]{/home/claude/optical-metrology-fitter/scripts/exp01_forward_model.py}
\subsection{\texttt{scripts/exp02\_inverse\_ambiguity.py} — 실험 2 — 역문제와 축퇴 (223줄)}
\lstinputlisting[firstnumber=1]{/home/claude/optical-metrology-fitter/scripts/exp02_inverse_ambiguity.py}
\subsection{\texttt{scripts/exp03\_precision\_crlb.py} — 실험 3 — CRLB 대 몬테카를로 (241줄)}
\lstinputlisting[firstnumber=1]{/home/claude/optical-metrology-fitter/scripts/exp03_precision_crlb.py}
\subsection{\texttt{scripts/exp04\_systematics.py} — 실험 4 — 계통오차 예산 (312줄)}
\lstinputlisting[firstnumber=1]{/home/claude/optical-metrology-fitter/scripts/exp04_systematics.py}
\subsection{\texttt{scripts/\_style.py} — 그림 스타일 (59줄)}
\lstinputlisting[firstnumber=1]{/home/claude/optical-metrology-fitter/scripts/_style.py}
\subsection{\texttt{run\_all.py} — 전체 실행기 (52줄)}
\lstinputlisting[firstnumber=1]{/home/claude/optical-metrology-fitter/run_all.py}

\section{맺음말 — 이 문서에서 반복되는 하나의 주제}

물리·코드·검증·실험을 관통하는 주제는 하나입니다.

\begin{center}
\emph{무엇을 실제로 측정했는가, 그리고 그것이 무엇을 결정하는가.}
\end{center}

\begin{itemize}[leftmargin=14pt,itemsep=3pt]
\item 프린지가 재는 것은 두께가 아니라 \textbf{광학 두께 $n\cdot d$}입니다 $\to$ $n$을 모르면 $d$를 모릅니다.
\item 단일 파장 반사율이 재는 것은 \textbf{$\cos2\delta$ 하나}입니다 $\to$ 차수와 측면이 축퇴합니다.
\item CRLB가 제한하는 것은 \textbf{분산}입니다 $\to$ 편향은 따로 재야 합니다.
\item 거칠은 막에서 ``두께''는 \textbf{규약}입니다 $\to$ 완벽한 두 장비가 불일치합니다.
\end{itemize}

이 프로젝트가 방어할 수 있는 진술은 다음 하나입니다.

\note{$\sigma_R=2\times10^{-4}$의 분광 반사계로 217\,nm SiO$_2$/Si를 측정할 때, 적합기는
통계적으로 효율적이며(몬테카를로 산포가 세 자릿수 잡음 범위에서 Cramér--Rao 하한의 2\%
이내) 랜덤 불확도는 4\,pm입니다. 현실적 정확도는 약 1.3\,nm이며, 모델 형태 —
모델링되지 않은 계면층, 거칠기 규약, 막 굴절률 가정 — 가 지배합니다.
\textbf{결과를 개선하려면 더 나은 옵티마이저가 아니라 더 나은 시료 지식이 필요합니다.}}

\vspace{6pt}
그리고 가장 큰 한계를 다시 명시합니다. \textbf{이 문서의 모든 숫자는 합성 데이터입니다.}
순방향 모델은 독립 구현체(Byrnes \texttt{tmm})와 $10^{-14}$ 수준으로 대조했지만,
\emph{역문제} 파이프라인은 실제 시료를 만난 적이 없습니다. 이를 검증하려면
기준 교정된 산화막 웨이퍼가 필요합니다.

\end{document}
